% Зөвхөн бүтэн ноорогт. 2026-11-28: хэрэглэгчийн шийдвэрээр багшид үзүүлэх ноорогоос гаргав (хэмжсэн үр дүн биш —
% хэрэгжилтийн төлөв ба хөгжүүлэлтийн асуултын тоо). Эцсийн үнэлгээний дараа «Системийн үнэлгээ» болж буцна.
\section{Системийн хэрэгжилтийн байдал}
\label{sec:res-system}

\ref{sec:system}-р хэсэгт тайлбарласан системийн хэсэг тус бүрийн 2026 оны 11-р сарын 27-ны байдлаарх төлөвийг \tref{tab:system_status}-д үзүүлэв.

\begin{table}[htbp]
\centering
\caption{Системийн хэсгүүдийн төлөв (2026.11.27)}
\label{tab:system_status}
\begin{tabular}{|p{6.2cm}|p{7.6cm}|}
\hline
\textbf{Хэсэг} & \textbf{Төлөв} \\
\hline
Бичгээр асууж хариулт авах вэб хуудас & Ажиллаж байна \\
\hline
Өрөө, анги, багш, хуваарийн асуултад сургуулийн хүснэгтээс хариулах & Ажиллаж байна \\
\hline
Олон асуудаг сэдэвт хадгалсан хариулт өгөх & Ажиллаж байна \\
\hline
Вэб сайтын хайлт (BM25, E5, \acs{RRF}) & Ажиллаж байна \\
\hline
Хариулт олдоогүй үед хандах алба, өрөөг заах & Ажиллаж байна; сэдэв ба албаны холбоосыг сургуулиар баталгаажуулна \\
\hline
Хариулт аль хүрээ, дэд сэдвээс гарсныг мэдлэгийн модоор харуулах & Ажиллаж байна \\
\hline
Мэдлэгийн модоор асуултын хүрээг тодорхойлж нэмэлт оноо өгөх & Хэрэгжүүлсэн; тохируулгын асуултаар хэмжиж идэвхжүүлэх эсэхийг шийднэ \\
\hline
Дуугаар асуух (Gemini 3.5 Transcribe) & Ажиллаж байна; таваас доошгүй хүний бичлэгээр үнэлнэ \\
\hline
Хариултыг дуугаар унших (edge-tts) & Ажиллаж байна; нэр, тоо, оныг сонсогч ойлгосныг үнэлнэ \\
\hline
Мэдээллийг хувилбартайгаар шинэчлэх, өмнөх хувилбар руу буцаах & Ажиллаж байна \\
\hline
Ресепшний дэлгэцэнд ажиллуулах & Эцсийн үнэлгээний өмнө суулгаж туршина \\
\hline
\end{tabular}
\end{table}

Системээс 2026 оны 11-р сарын 27-нд шууд хуулсан гурван хариултыг доор үзүүлэв. Хүний нэр агуулсан хэсгийг «…» тэмдгээр орхив.
\begin{itemize}
    \item \textbf{Вэб сайтаас хариулсан:} «Сургууль хэдэн онд байгуулагдсан бэ?» → «Шинэ Монгол технологийн коллеж 2014 оны 9-р сарын 15-нд үүсгэн байгуулагдсан; анхны 93 оюутан 3 бүлэгтэй хичээллэж эхэлсэн. …»
    \item \textbf{Сургуулийн хүснэгтээс хариулсан:} «Номын сан хэддүгээр давхарт байдаг вэ?» → «Номын сан: 10-р давхар, 1006 тоот. Ажилтнууд: …»
    \item \textbf{Хариулт олдоогүй үед чиглүүлсэн:} «Хүүхдээ давтлагад оруулмаар байна» → «Энэ талаар мэдээлэл олдсонгүй. Лавлах газар: Сургалтын алба (4-р давхар, 406 тоот, ажлын цаг 08:00–17:00). Сургуулийн ерөнхий утас +976 75777799, имэйл info@nmct.edu.mn.»
\end{itemize}

Хөгжүүлэлтийн явцад системийг энэ ажлын хүрээнд бичсэн асуултуудаар шалгаж алдааг засварласан боловч засварт ашигласан асуултад зөв хариулсан тоо системийн чанарыг бодитоор илэрхийлэхгүй тул энд тайлагнаагүй бөгөөд эцсийн үр дүнг \ref{sec:evaluation-method}-р хэсэгт тодорхойлсон бодит асуултаар гаргана.

Хариулт аль замаар гарч буйг болон мэдлэгийн модны нэмэлт онооны нөлөөг хөгжүүлэлтийн асуултаар шалгав. Хөгжүүлэлтийн 40 асуултын 18-д хадгалсан хариултаас, 10-т сургуулийн хүснэгтээс хариулж, 12-т хариулаагүй («олдсонгүй», ихэвчлэн хандах газрын хамт) бөгөөд вэб сайтын хайлтын хэсгээр шууд хариулсан асуулт байгаагүй (2026.11.27). Иймд одоогийн системд хайлтын үүрэг нь хариултыг шууд гаргахаас илүүтэй эх сурвалжийн хуудсыг олоход оршиж байна. Мэдлэгийн модны нэмэлт оноог унтраасан ба асаасан үед зөв хуудас хайлтын эхний үр дүнд орсон асуултын тоог \tref{tab:tree_boost}-д харуулав.

\begin{table}[htbp]
\centering
\caption{Мэдлэгийн модны нэмэлт оноо: унтраасан ба асаасан (хөгжүүлэлтийн 76 асуулт, 2026.11.27)}
\label{tab:tree_boost}
\begin{tabular}{|p{6.4cm}|r|r|}
\hline
\textbf{Зөв хуудас} & \textbf{Унтраасан} & \textbf{Асаасан} \\
\hline
Эхний 1 үр дүнд & 55 & 56 \\
\hline
Эхний 3 үр дүнд & 66 & 66 \\
\hline
Эхний 10 үр дүнд & 72 & 72 \\
\hline
\end{tabular}
\end{table}

Нэмэлт оноо 76 асуултын есийнх нь зөв хуудасны байрыг өөрчилсний зургаа нь сайжирч, гурав нь муудсан. Иймд мод одоогоор хайлтын чанарт бус, мэдээллийн хамрах хүрээг тодорхойлох, ангилах, дутууг илрүүлэхэд хэрэгтэй байна. Эдгээр асуултыг энэ ажлын хүрээнд бичиж, зарим нь хөгжүүлэлтэд ашиглагдсан тул тоо нь эцсийн дүн биш юм.

